\section{An exact weight-five normalization}
\label{sec:weight-five-normalization}

Let $G=\widehat{\mathrm{GT}}$, let $H$ be its Galois image, and write
the radical and cubic normalizations as
\[
 D=N_{\mathrm{rad}},\qquad
 D_3=D\cap\ker m_3,\qquad J_3=D_3\cap H.
\]
The established factorizations give $G=D_3H$.  All Magnus coefficients
in this section are $2$-adic.  For $g=(1,f)\in D_3$, the pro-$2$
Magnus expansion of $f$ starts in degree five: the degree-three and
degree-four Hall coordinates are respectively
$(m_3,-m_3)$ and $(-m_3,m_3,-m_3)$ by the preceding low-weight
calculation.  Define
\[
 m_5(g)=[S^4T]f,\qquad x=1+S,\quad y=1+T.
\]

\begin{theorem}\label{thm:weight-five-exact-image}
The exact integral images are
\[
 m_5(D_3)=m_5(J_3)=32\mathbf Z_2.
\]
In particular,
\[
 \eta_5:D_3\longrightarrow\mathbf F_2,
 \qquad \eta_5(g)=m_5(g)/32\pmod2,
\]
is a continuous $G$-invariant character with nonzero arithmetic
restriction.  The kernel of $m_5$,
\[
 D_5=D_3\cap\ker m_5,
 \qquad J_5=D_5\cap H=J_3\cap\ker m_5,
\]
is a proper closed normal subgroup satisfying
\[
 D_3=D_5J_3,\qquad G=D_5H.
\]
The relative abelian quotient refines canonically to
\[
 B\simeq D_3/\overline{[G,D_3]J_3}
       \simeq D_5/\overline{[G,D_5]J_5}.
\]
\end{theorem}

The exact image is the case $r=2$ of
Theorem~\ref{thm:uniform-odd-arithmetic-image}. We give its explicit
arithmetic value and the resulting extension obstruction.

\begin{proposition}\label{prop:weight-five-upper-bound}
For $g\in D_3$, one has $m_5(g)\in32\mathbf Z_2$.
\end{proposition}

\begin{proof}
Since $D_3=C_1$, apply Theorem~\ref{thm:uniform-odd-arithmetic-image}
with $r=2$. Equivalently, the integral measure has zero mass and zero
second moment, so the integer-valued polynomial of degree four in
Lemma~\ref{lem:uniform-odd-coordinate-bound} gives the factor $2^5$.
\end{proof}

\subsection{An arithmetic generator of the exact image}

The arithmetic lower bound uses a weight-five homomorphism on
the entire multiplier-one subgroup.  For such an element, write
\[
 a_3=[S^2T]f,\qquad b_5=[S^4T]f,
 \qquad \tau_5(f)=12b_5-11a_3\in\mathbf Z_2.
\]

\begin{lemma}\label{lem:tau-five}
The map $\tau_5$ is a continuous homomorphism on $\ker\lambda$,
and
\[
 \tau_5(hzh^{-1})=\lambda_2(h)^5\tau_5(z)
 \qquad(h\in G,\ z\in\ker\lambda).
\]
On $D_3$, it is equal to $12m_5$.
\end{lemma}

\begin{proof}
Let $R(X,Y)$ be the Alexander coefficient of $f$ on $[x,y]$,
and let
\[
 B_f(X,Y)=1+(X-1)(Y-1)R(X,Y).
\]
The action on the rank-one Alexander module is semilinear.
Consequently, on the multiplier-one subgroup the cofactors
multiply, and conjugation by multiplier $t$ sends $B_f(X,Y)$
to $B_f(X^t,Y^t)$.  It follows that
\[
 [u^4v]\log B_f(e^u,e^v)
\]
is an additive weight-five $\mathbf Q_2$-valued homomorphism.
Only terms linear in $Y-1$ contribute to this coefficient, so
the logarithm can be omitted for this extraction.  The boundary
Magnus coefficients and the earlier identity $[S^3T]f=-a_3$ give
\[
 (X-1)R(X,1)
 =a_3(X-1)^2-a_3(X-1)^3+b_5(X-1)^4+O((X-1)^5).
\]
Now
\[
 [u^4](e^u-1)^2=\frac7{12},\quad
 [u^4](e^u-1)^3=\frac32,\quad
 [u^4](e^u-1)^4=1.
\]
Thus the extracted homomorphism is $b_5-11a_3/12$.
Multiplication by $12$ proves the asserted integral formula
and its exact transformation rule.  The last assertion uses
$a_3=0$ on $D_3$.
\end{proof}

\begin{proposition}\label{prop:weight-five-arithmetic-generator}
There exists $w\in J_3$ with $m_5(w)\in32\mathbf Z_2^\times$.
\end{proposition}

\begin{proof}
Simultaneous rational Kummer surjectivity over
$\mathbf Q^{\mathrm{ab}}$ supplies $\sigma\in H\cap\ker\lambda$
with
\[
 \rho_2^I(\sigma)=2\in\widehat{\mathbf Z},\qquad
 \rho_q^I(\sigma)=0\quad(q\text{ odd}).
\]
The established boundary congruence
$\rho_2^I\equiv-2m_3\pmod4$ makes $a_3(\sigma)$ odd.
Therefore $\tau_5(\sigma)=12b_5(\sigma)-11a_3(\sigma)$ is odd.

Choose $h\in H$ with $\lambda_2(h)=3$ and put
$t=\lambda(h)\in\widehat{\mathbf Z}^{\times}$.  For
$c\in\widehat{\mathbf Z}$ and $z\in H\cap\ker\lambda$, define
\[
 \mathcal D_c(z)=hzh^{-1}z^{-c}.
\]
The power $z^{-c}$ is the standard continuous profinite power
in the procyclic closure of $z$.

Every rational Kummer coordinate of $\mathcal D_t(\sigma)$
is zero, because its conjugation weight is $t$.  Set
\[
 w=\mathcal D_{t^3}\bigl(\mathcal D_t(\sigma)\bigr).
\]
All rational Kummer coordinates remain zero.  Its cubic
coordinate also vanishes, since conjugation on that coordinate
has weight $t_2^3$.  Hence $w\in J_3$.

Lemma~\ref{lem:tau-five} and the continuity of powers give the
exact equality
\[
 \tau_5(w)=(3^5-3^3)(3^5-3)\tau_5(\sigma)
          =51840\tau_5(\sigma).
\]
This equality holds for the exact homomorphism $\tau_5$ with the
stated conjugation weight.  Since $w\in D_3$,
\[
 m_5(w)=\frac{51840}{12}\tau_5(\sigma)
       =4320\tau_5(\sigma)
       =32\cdot135\,\tau_5(\sigma),
\]
which is a generator of $32\mathbf Z_2$.
\end{proof}

\begin{proof}[Proof of Theorem~\ref{thm:weight-five-exact-image}]
On $D_3$, $m_5$ is a continuous additive homomorphism, with
conjugation weight $\lambda_2^5$; this follows either from its
leading Magnus degree or from Lemma~\ref{lem:tau-five} and
the torsion-freeness of $\mathbf Z_2$.  The upper bound and
the arithmetic generator show that its compact images on
both $D_3$ and $J_3$ are exactly $32\mathbf Z_2$.  Dividing
by $32$ and reducing modulo two gives a $G$-invariant
character, since every $2$-adic unit has odd fifth power.

Equality of the images gives $D_3=D_5J_3$.  Normality follows
from the scalar conjugation rule, properness from the
arithmetic generator, and $G=D_5H$ follows from $G=D_3H$.
For the quotient assertion, if $G=EH$ with $E$ closed normal,
then
\[
 E/\overline{[G,E](E\cap H)}
 \simeq G/\overline{[G,E]H}.
\]
If $E=E'(E\cap H)$ and $G=E'H$, the closures
$\overline{[G,E]H}$ and $\overline{[G,E']H}$ coincide: in
the quotient by the latter, $H$ is trivial and the image
of $E'$ is central.  Apply this with $E=D_3$ and $E'=D_5$.
\end{proof}

If $J_3=G_{L_3}$ and $L_5=\overline{\mathbf Q}^{\,J_5}$,
then $\operatorname{Gal}(L_5/L_3)\simeq\mathbf Z_2$.
The character $\eta_5$ has nonzero arithmetic restriction on $D_3$.
The same witness shows precisely why it does not give another invariant
character on the original normalized subgroup.

\begin{corollary}\label{cor:weight-five-nonextension}
The character $\eta_5$ does not extend to a $G$-invariant continuous
$\mathbf F_2$-character of $D=N_{\mathrm{rad}}$, and hence does not
extend to one of $N$.  Nevertheless $D_5$ is a strictly smaller
normalization with the same relative abelian quotient.  The kernel of
restriction from $X_2$ to $H^1(D_5,\mathbf F_2)$ is exactly
\[
 \bigoplus_{q\text{ odd prime}}\mathbf F_2\psi_q
       \oplus\mathbf F_2\eta_3.
\]
The kernel of restriction from $H^1(D,\mathbf F_2)^G$ to
$H^1(D_5,\mathbf F_2)$ is exactly $\mathbf F_2\eta_3|_D$.
For every odd prime $p$, restriction from $H^1(N,\mathbf F_p)^G$
to $H^1(D_5,\mathbf F_p)$ is injective.
\end{corollary}

\begin{proof}
Use the elements $h,\sigma,t$ of
Proposition~\ref{prop:weight-five-arithmetic-generator}, and put
$z=\mathcal D_t(\sigma)\in D\cap H$.  Its witness is
\[
 w=hzh^{-1}z^{-t^3}=[h,z]z^{1-t^3}\in D_3\cap H,
 \qquad \eta_5(w)=1.
\]
Since $t$ is a profinite unit, $1-t^3\in2\widehat{\mathbf Z}$.
The continuous profinite power $z^{(1-t^3)/2}$ is defined, so
$w$ belongs to the closed subgroup generated by $[G,D]$ and the
squares in $D$.  Every invariant quadratic character of $D$ kills
that subgroup.  It therefore takes the value zero at $w$, which is
incompatible with extending $\eta_5$.  Restricting an invariant
character of $N$ to $D$ proves the second nonextension assertion.

If $\chi\in X_2$ vanishes on $D_5$, its restriction to $D_3$
factors through $D_3/D_5\simeq\mathbf Z_2$.  That restriction is
either zero or $\eta_5$.  The latter possibility has just been
excluded, so $\chi$ vanishes on $D_3$.  Conversely every character
vanishing on $D_3$ vanishes on $D_5$.  Corollary~\ref{cor:cubic-radical-joint-image}
identifies this common kernel as the displayed direct sum.
The same argument on $D$, using $D/D_3\simeq8\mathbf Z_2$,
gives the one-dimensional kernel there.

For odd $p$, a character of $D_3$ vanishing on $D_5$ must be zero,
because $D_3/D_5$ is pro-$2$.  A character of $N$ vanishing on
$D_3$ factors through the joint rational Kummer and cubic quotient.
The cubic factor is pro-$2$, and conjugation by multiplier $-1$
acts by inversion on every rational Kummer coordinate.  Invariance
therefore kills the remaining factor in $\mathbf F_p$.
\end{proof}
